\section{Resplendence, Definability, and Atomic Compactness}
\label{sec:p03-resplendence-atomic-compactness}

An expansion interprets new symbols on the existing domain. Compactness
with an elementary diagram initially supplies an elementary extension,
not such an expansion. The witness-selection construction below is the
step which brings the interpretations back to the fixed model. Its
countable recursive-language presentation follows the actual Henkin
argument of \cite[\S1.6, Definition~1.22 and Theorem~1.23, pp.~17--18]
{Engstrom2004Expansions}, with a finite-tuple version and an explicit
recursive-saturation schedule supplied here. Atomic compactness has a
different parameter and variable scope; we reconstruct the diagram and
retraction argument in \cite[\S2, pp.~100--102]{Delhomme1999AtomicGraphs}.
The recursive-language expansion presentation and the atomic-system
presentation retain their different language, parameter and variable scopes.

\subsection{Effective syntax and fixed-model expansion}

\begin{hbconvention}{def:p03-resplendence-coding}{External recursive syntax, not a computable model}
For the effective assertions let \(L\) be a countable, one-sorted,
finitary first-order language with nonempty domains, whose symbol codes,
arities and term/formula syntax are recursive subsets of \(\mathbb N\).
Fix a recursive coding of finite strings and a recursively checkable sound
and complete first-order proof calculus with equality. A recursive extension
\(L^+\) has the same sort and effectively coded new symbols; it may be
countable. Parameter tuples are finite. Their additional constant codes
are formal names, not a computable numbering of the domain of \(M\).
Apart from the expressly shared names for \(\bar a\), the names later
adjoined for members of \(M\) lie in a tagged namespace disjoint from
all new \(L^+\)-symbols. A displayed \(c_k\) in the named-model example
is a base-language symbol, not one of these additional parameter names.
Recursive and computably enumerable refer to these external standard
natural-number codes. Neither \(\operatorname{Th}(M)\), the elementary
diagram, the presentation of \(M\), nor the choices in the construction
are assumed decidable. We work in \(\mathsf{ZF}+\mathsf{AC}\), using
\cref{thm:fol-completeness,thm:fol-compactness,thm:downward-ls,thm:elementary-chain}
at their stated scope. The atomic-compactness results later require no
recursive-language hypothesis.
\end{hbconvention}

\begin{hbdefinition}{def:p03-resplendence-recursive-saturation}{Recursive saturation over finite parameters}
The model \(M\) is recursively saturated if every recursive set
\(p(\bar x,\bar a)\) of formulas in a fixed finite variable block,
with a fixed finite tuple \(\bar a\) from \(M\), which is finitely
satisfiable in \(M\), has a simultaneous realization in \(M\).
Membership means membership of the syntactic formulas before the formal
parameter names are interpreted. Types need not be complete.
\end{hbdefinition}

\begin{hblemma}{thm:p03-resplendence-padding}{proved=annex:p03-resplendence-padding}{Computably enumerable types can be made recursive}
Every computably enumerable type over a fixed finite parameter tuple has
a recursive type with exactly the same realizations, in the same finite
variable block. Thus recursive saturation realizes the former as well.
\end{hblemma}

\begin{hbdefinition}{def:p03-resplendence-expansion}{Two declared expansion properties}
The recursive-theory resplendence convention used in
\cite[Definition~1.22]{Engstrom2004Expansions} is this: for every finite
\(\bar a\in M\), recursive extension \(L^+\supseteq L(\bar a)\), and
recursive \(L^+\)-theory \(T\) consistent with
\(\operatorname{Th}_{L(\bar a)}(M,\bar a)\), an \(L^+\)-expansion
\(M^+\) of \((M,\bar a)\) satisfies \(T\).
We also explicitly distinguish finite-sentence resplendence: the same
requirement for a sentence in a language with finitely many new symbols.
The theorem below proves the recursive-theory property, and therefore
the finite-sentence property, from a separate realization hypothesis.
The definition itself is not an existence proof.
\end{hbdefinition}

\begin{hblemma}{thm:p03-resplendence-diagram}{proved=annex:p03-resplendence-diagram}{The finite-parameter theory and the full elementary diagram}
Let \(S\) be any expanded-language theory mentioning only finitely many
names from \(M\), collected in \(\bar a\). Then
\(S\cup\operatorname{Th}(M,\bar a)\) is consistent if and only if
\(S\cup\operatorname{ElDiag}(M)\) is consistent. A model of the second
theory contains an elementary copy of \(M\); it need not have that copy
as its whole domain.
\end{hblemma}

\begin{hblemma}{thm:p03-resplendence-witness}{proved=annex:p03-resplendence-witness}{Actual witnesses from a recursively saturated model}
Let \(M\) be recursively saturated, \(S\) a computably enumerable
\(L^+\)-theory with finite parameter footprint in \(M\), and
\(p(\bar x,\bar a)\) a computably enumerable expanded-language type
with finite parameter footprint. If
\[
 \operatorname{ElDiag}(M)\cup S\cup p(\bar d,\bar a)
\]
is consistent for fresh constants \(\bar d\), then there is an actual
tuple \(\bar b\in M\) such that
\(\operatorname{ElDiag}(M)\cup S\cup p(\bar b,\bar a)\) is consistent.
This includes the one-formula Henkin witness case; it does not assume
that \(M\) is fully saturated.
\end{hblemma}

\begin{hbtheorem}{thm:p03-resplendence-countable}{proved=annex:p03-resplendence-countable}{Countable recursive saturation supplies resplendence}
For countable \(M\) in the declared recursive language, recursive
saturation is equivalent to recursive-theory resplendence. More strongly,
if \(M\) is recursively saturated and \(T\) is a computably enumerable
theory in a recursive extension \(L^+\supseteq L(\bar a)\), with
\(\operatorname{Th}(M,\bar a)\cup T\) consistent, there is a
\emph{recursively saturated} expansion \(M^+\models T\) on exactly
the domain of \(M\), with the old symbols and \(\bar a\) unchanged.
The annex schedules both Henkin witnesses and all recursive expanded
types; it does not delegate that last assertion to a satisfaction class.
\end{hbtheorem}

\begin{hbtheorem}{thm:p03-resplendence-extension}{proved=annex:p03-resplendence-extension}{Building countable recursively saturated elementary extensions}
Every countable model in a recursive language has a countable recursively
saturated elementary extension. Combined with the preceding theorem this
gives a recursively resplendent elementary extension. The union-of-chains
construction is explicit; it does not assert that the original model was
already resplendent or that an arbitrary prescribed expansion is recursively
saturated. This supplies the countable case of
\cite[Proposition~1.4, p.~10]{Engstrom2004Expansions}.
\end{hbtheorem}

\subsection{Homogeneity and a parameter-sensitive definability criterion}

\begin{hbtheorem}{thm:p03-resplendence-homogeneity}{proved=annex:p03-resplendence-homogeneity}{Finite elementary maps extend, and countable back-and-forth}
In a recursively saturated model, if finite tuples \(\bar a,\bar b\)
have the same complete \(L\)-type, then for each \(c\in M\) there is
\(d\in M\) with
\(\operatorname{tp}(\bar a,c)=\operatorname{tp}(\bar b,d)\).
For countable \(M\), every finite elementary self-map extends to an
automorphism. The proof uses recursive comparisons of formulas with
the actual tuple \((\bar a,\bar b,c)\), not a claim that its complete
type is recursive \cite[Proposition~1.6, pp.~10--11]{Engstrom2004Expansions}.
\end{hbtheorem}

\begin{hbtheorem}{thm:p03-resplendence-definability}{proved=annex:p03-resplendence-definability}{Definability when the relation expansion is recursively saturated}
Let \(M\) be countable in recursive \(L\), let \(R\subseteq M^n\),
where \(n\) is finite, and assume the \emph{given} expansion \((M,R)\)
is recursively saturated. For a finite parameter tuple \(\bar a\),
the relation \(R\) is definable by an \(L(\bar a)\)-formula if and
only if every automorphism of the \(L\)-reduct fixing \(\bar a\)
preserves \(R\). The annex constructs the finite Boolean combination
of base formulas, or else realizes two equal-base-type tuples with
different \(R\)-membership and constructs a separating automorphism.
This is the declared local criterion; it is not an assertion that
arbitrary automorphism-invariant subsets are definable.
\end{hbtheorem}

\subsection{Atomic systems, purity, and actual retractions}

\begin{hbdefinition}{def:p03-resplendence-atomic}{Atomic and positive compactness}
Here \(L\) may be any set-sized finitary one-sorted language.
Let \(A\) be a nonempty \(L\)-structure, let \(t_1,t_2\) be
\(L(A)\)-terms of its object sort, and let \(C,D\) denote structures.
In Delhomme's native atomic presentation, \(t_1\equiv t_2\) is a
logical identity atom: its interpretation compares those two terms'
object values in their sort.
An atom is either \(t_1\equiv t_2\) or \(R(t_1,\ldots,t_k)\).
This term/object use of \(\equiv\) is distinct from
\(C\equiv D\) with \emph{structures} as operands, which denotes
elementary equivalence. Ordinary metamathematical \(=\) between values,
maps or sets retains its existing meaning. The native identity atom is
preserved in this presentation and its homomorphism diagram; its comparison
with the host logical atom \(t_1=t_2\) is proved below.
A system
\(\Phi(X)\) is a set of formulas with parameters from \emph{all} of
\(A\), and free variables from an arbitrary set \(X\). Each formula
has finitely many variables, but \(X\) and the parameter set need not be
finite or countable. The system is finitely solvable if every finite
subsystem has an assignment in \(A\), and solvable if one assignment
in \(A^X\) satisfies all its members. Atomic compactness requires this
implication for atomic systems. Positive compactness requires it for
systems built from atomic formulas by \(\land,\lor,\exists,\forall\),
without negation. Positive primitive (pp) formulas use \(\land,\exists\)
only. These are the scopes of \cite[\S2]{Delhomme1999AtomicGraphs}.
Atomic compactness is not atomicity of a model (isolation of its types).
\end{hbdefinition}

\begin{hbdefinition}{def:p03-resplendence-purity}{Pure extensions and retractions}
An inclusion \(A\subseteq B\) is pure if every pp formula with a finite
tuple of parameters from \(A\) true in \(B\) is true in \(A\).
A retraction is an \(L\)-homomorphism \(r:B\to A\) fixing \(A\)
pointwise. Homomorphisms commute with functions and preserve relations;
they need not reflect relations, be injective, or be elementary.
Every elementary extension is pure. Our retractions are onto because
they fix \(A\); surjectivity matters for universal positive formulas.
\end{hbdefinition}

\begin{hblemma}{thm:p03-resplendence-positive-maps}{proved=annex:p03-resplendence-positive-maps}{Positive formulas under homomorphisms}
Every homomorphism preserves positive existential formulas. Every onto
homomorphism preserves all positive formulas. In particular a retraction
returns a solution of any positive system with parameters in \(A\)
from its extension to \(A\).
For any structure \(C\), assignment \(s\), and same-sort terms,
the native and host identity atoms have the satisfaction bridge
\[
 C,s\models(t_1\equiv t_2)
 \quad\Longleftrightarrow\quad t_1^C(s)=t_2^C(s)
 \quad\Longleftrightarrow\quad C,s\models(t_1=t_2).
\]
The middle equality is identity of object values in the metatheory;
the outside expressions are satisfaction of two declared logical atoms.
\end{hblemma}

\begin{hbtheorem}{thm:p03-resplendence-atomic-retractions}{proved=annex:p03-resplendence-atomic-retractions}{The diagram construction for atomic compactness}
For any \(A\), the following are equivalent: atomic compactness;
being a retract of every pure extension; being a retract of every
elementary extension; and positive compactness. The construction introduces
one unknown for each element of an extension and writes every function,
relation and fixed-parameter constraint. Solving that atomic system
produces the retraction. Conversely compactness first realizes a finitely
solvable system in an elementary extension, and the actual retraction
returns it to \(A\) \cite[\S2, p.~101]{Delhomme1999AtomicGraphs}.
\end{hbtheorem}

\begin{hbtheorem}{thm:p03-resplendence-topological}{proved=annex:p03-resplendence-topological}{Compact topological structures give atomic solutions}
If \(A\) has a compact Hausdorff topology, every basic relation is closed
and every function is continuous, then \(A\) is atomic compact.
This includes finite structures with discrete topology and the Boolean
algebra \(\{0,1\}^{\mathbb N}\) with coordinatewise operations and
product topology. The proof works for arbitrary set-sized unknowns, not
just finitely many variables \cite[\S2, pp.~101--102]{Delhomme1999AtomicGraphs}.
\end{hbtheorem}

\begin{hbdefinition}{def:p03-resplendence-reduced-power}{The filter quotient used by the atomic argument}
Let \(I\ne\varnothing\) be a set and \(F\) a proper filter on
\(\mathcal P(I)\). For \(u,v\in A^I\), put
\(u\sim_F v\) if \(\{i:u(i)=v(i)\}\in F\). Interpret functions
coordinatewise on equivalence classes, and interpret a basic relation by
\[
 A^I/F\models R([u_1],\ldots,[u_k])
 \quad\Longleftrightarrow\quad
 \{i:A\models R(u_1(i),\ldots,u_k(i))\}\in F.
\]
This includes nullary relations. The diagonal sends \(a\) to its constant
function class. An ultrapower uses an ultrafilter. The proof below verifies
the quotient and the required pp truth calculation before using purity.
\end{hbdefinition}

\begin{hbtheorem}{thm:p03-resplendence-atomic-powers}{proved=annex:p03-resplendence-atomic-powers}{Reduced-power and ultrapower witnesses}
Atomic compactness is also equivalent to \(A\) being a retract of every
diagonally embedded reduced power \(A^I/F\), for proper filters \(F\),
and equivalent to the analogous assertion for every ultrapower.
The annex proves the pp reduction and constructs a solution from finite
subsystem witnesses and the cone filter; it does not invoke a false
full first-order \L o\'s theorem for arbitrary reduced products.
\end{hbtheorem}

\subsection{Worked expansions and separating failures}

\begin{hbtheorem}{thm:p03-resplendence-examples}{proved=annex:p03-resplendence-examples}{Fixed domains and the two different compactness demands}
The following examples have the stated full constructions.
\begin{enumerate}
\item The equality-only model on \(\mathbb N\) is recursively saturated.
On that same domain the new function \(f\) and predicate \(P\), given
by \(f(2k)=2k+1\), \(f(2k+1)=2k\), and \(P=2\mathbb N\), satisfy
\(f^2=\mathrm{id}\), \(f(x)\ne x\), and
\(P(f(x))\leftrightarrow\neg P(x)\). This is a fixed-model expansion,
not an elementary extension. Neither \(P\) nor the graph of \(f\)
is definable in the equality reduct over a finite parameter tuple.
\item In the recursive language with a constant \(c_k\) for every
\(k\in\mathbb N\), interpret \(c_k=k\) on \(\mathbb N\).
The recursive type \(\{x\ne c_k:k\in\mathbb N\}\) is finitely
satisfiable but omitted. Adding a new constant to realize it is consistent
with the complete theory, yet impossible on this fixed named domain.
This model is nevertheless atomic compact. Every subset is invariant
under its automorphisms; an infinite coinfinite subset is not definable
over finite parameters. Invariance without the expansion hypothesis
of the definability theorem is insufficient.
\item A countable recursively saturated elementary extension of
\((\mathbb N,0,S,<)\) exists by the actual chain construction above.
It is recursively resplendent, but not atomic compact: the atomic
system \(\{a<x:a\in M\}\) is finitely solvable and has no solution.
Allowing all parameters is the decisive difference.
\end{enumerate}
\end{hbtheorem}

\begin{proofboundary}{def:p03-resplendence-original-boundary}{Countability and limits of the constructions}
The countable recursive construction, parameter-sensitive criterion, and
arbitrary atomic systems have their distinct hypotheses. Countability is
used to schedule all named-domain requirements; no uncountable recursive
saturation-to-resplendence assertion is made. The recursive-theory theorem
does not assert expandability for arbitrary theories without its coding
and finite-parameter conditions. The retraction characterization assumes
atomic compactness of the base; it does not construct a pure-injective
hull or establish a hull's minimality or uniqueness. The sufficient compact
topology condition is not asserted necessary for atomic compactness.
\end{proofboundary}
